\markdownRendererDocumentBegin
\markdownRendererSectionBegin
\markdownRendererHeadingOne{Rapporto di validazione Allen–Cahn in OOFEM}\markdownRendererInterblockSeparator
{}\markdownRendererStrongEmphasis{Revisione scientifica indipendente degli artefatti disponibili al 26 ottobre 2026.} Il rapporto sintetizza i test non forzati/non costruiti, i confronti con soluzioni analitiche e riferimenti pubblicati, i riferimenti semidiscreti indipendenti e le prove energetiche già presenti. Non sono state eseguite nuove simulazioni per redigere questo documento.\markdownRendererInterblockSeparator
{}La classificazione per componente è nella \markdownRendererLink{matrice finale}{validation\markdownRendererUnderscore{}matrix\markdownRendererUnderscore{}final.md}{validation_matrix_final.md}{} e in \markdownRendererLink{CSV}{validation\markdownRendererUnderscore{}matrix\markdownRendererUnderscore{}final.csv}{validation_matrix_final.csv}{}; l’audit quantitativo è in \markdownRendererLink{validation\markdownRendererUnderscore{}metrics\markdownRendererUnderscore{}audit.md}{validation\markdownRendererUnderscore{}metrics\markdownRendererUnderscore{}audit.md}{validation_metrics_audit.md}{}, l’inventario file in \markdownRendererLink{validation\markdownRendererUnderscore{}inventory.md}{validation\markdownRendererUnderscore{}inventory.md}{validation_inventory.md}{} e le tre ripetizioni operative in \markdownRendererLink{reproducibility\markdownRendererUnderscore{}runs.md}{reproducibility\markdownRendererUnderscore{}runs.md}{reproducibility_runs.md}{}.\markdownRendererInterblockSeparator
{}\markdownRendererSectionBegin
\markdownRendererHeadingTwo{1. Modelli implementati}\markdownRendererInterblockSeparator
{}Il materiale \markdownRendererCodeSpan{AllenCahnMaterial} usa flusso \markdownRendererInlineMath{-k\nabla u}, reazione bulk \markdownRendererInlineMath{m(u-u^3)} e reazione di linea \markdownRendererInlineMath{\beta m(\psi-\psi^3)}, con \markdownRendererInlineMath{\psi=u|_\Gamma}. Gli elementi \markdownRendererCodeSpan{line1ht} condividono i gradi di libertà della traccia con il bulk: non è un sistema con un secondo campo superficiale indipendente. I bilanci sotto assumono linee che ricoprono la frontiera dinamica una volta, coefficienti positivi, assenza di carichi/lavoro al bordo e regolarità sufficiente.\markdownRendererInterblockSeparator
{}\markdownRendererStrongEmphasis{Parabolico bulk-superficie.} La forma continua corrispondente è\markdownRendererInterblockSeparator
{}\markdownRendererDisplayMath{\begin{aligned} c u_t-k\Delta u &=m(u-u^3) &&\text{in }\Omega,\\ c\psi_t-k\Delta_\Gamma\psi+k\partial_nu &=\beta m(\psi-\psi^3) &&\text{su }\Gamma,\\ \psi&=u|_\Gamma. \end{aligned}}\markdownRendererInterblockSeparator
{}Il coefficiente \markdownRendererCodeSpan{c} è quello usato sui DOF bulk e di linea nei casi parabolici esaminati; la capacità superficiale indipendente non è parametrizzata separatamente. \markdownRendererCodeSpan{surfacebeta} pesa la sola reazione/potenziale superficiale. La forma debole bulk produce il termine normale che si accoppia e si cancella nel bilancio; non è un carico indipendente.\markdownRendererInterblockSeparator
{}\markdownRendererStrongEmphasis{Iperbolico OOFEM (\markdownRendererCodeSpan{wave 1}).} È inerziale soltanto il bulk; la superficie resta del primo ordine:\markdownRendererInterblockSeparator
{}\markdownRendererDisplayMath{\begin{aligned} \rho u_{tt}+\rho\zeta u_t-k\Delta u &=m(u-u^3) &&\text{in }\Omega,\\ c_\Gamma\psi_t-k\Delta_\Gamma\psi+k\partial_nu &=\beta m(\psi-\psi^3) &&\text{su }\Gamma,\\ \psi&=u|_\Gamma. \end{aligned}}\markdownRendererInterblockSeparator
{}\markdownRendererCodeSpan{rho} è la massa inerziale del bulk, \markdownRendererCodeSpan{wavedamping} \markdownRendererInlineMath{\zeta} moltiplica la massa bulk e \markdownRendererCodeSpan{surfacecapacity} \markdownRendererInlineMath{c_\Gamma} pesa la dinamica superficiale. Non è implementato/verificato un termine di inerzia superficiale \markdownRendererInlineMath{\psi_{tt}}. I test Newmark usano \markdownRendererInlineMath{\beta_N=1/4}, \markdownRendererInlineMath{\gamma_N=1/2}.\markdownRendererInterblockSeparator
{}
\markdownRendererSectionEnd \markdownRendererSectionBegin
\markdownRendererHeadingTwo{2. Metodologia di validazione}\markdownRendererInterblockSeparator
{}Le prove coprono tipi di riferimento diversi, da non confondere:\markdownRendererInterblockSeparator
{}\markdownRendererUlBeginTight
\markdownRendererUlItem \markdownRendererStrongEmphasis{Soluzione analitica del PDE bulk:} onda viaggiante iperbolica Nizovtseva, kink stazionario. Il dominio numerico è finito, quindi Neumann naturale approssima le code dell’onda.\markdownRendererUlItemEnd 
\markdownRendererUlItem \markdownRendererStrongEmphasis{Soluzione esatta del sistema FEM semidiscreto:} modo coseno lineare; isola l’integrazione temporale Newmark, non l’errore spaziale del PDE.\markdownRendererUlItemEnd 
\markdownRendererUlItem \markdownRendererStrongEmphasis{Riferimento ODE indipendente:} oscillatore Allen–Cahn uniforme e sistema accoppiato bulk-inerziale/superficie-parabolico. Il riferimento accoppiato assembla le matrici Q1/P1 in Python e integra con DOP853; è un controllo indipendente, ma la sua concordanza va verificata, non assunta.\markdownRendererUlItemEnd 
\markdownRendererUlItem \markdownRendererStrongEmphasis{Confronto bibliografico:} Okumura, energia digitalizzata dalla Fig. 4.4; Xiao–Chen, confronto visivo con \markdownRendererInlineMath{\varepsilon=0.1} assunto perché il valore non è dichiarato nel caso pubblicato.\markdownRendererUlItemEnd 
\markdownRendererUlItem \markdownRendererStrongEmphasis{Convergenza interna:} norme tra mesh o passi temporali differenti. La mesh/soluzione più fine è un riferimento numerico, non la soluzione esatta.\markdownRendererUlItemEnd 
\markdownRendererUlItem \markdownRendererStrongEmphasis{Bilancio energetico:} ricostruzione degli addendi e dissipazioni dai campi FEM salvati. È una verifica numerica a passo finito e dipende dal post-processing; non equivale a una prova dell’identità discreta del solver.\markdownRendererUlItemEnd 
\markdownRendererUlEndTight \markdownRendererInterblockSeparator
{}Le metriche sono prevalentemente RMS/max nodali oppure norme FEM/trapezoidali su griglie ricostruite; il tipo di norma è riportato con ciascun risultato. Le ripetizioni operative sono descritte in \markdownRendererLink{reproducibility\markdownRendererUnderscore{}runs.md}{reproducibility\markdownRendererUnderscore{}runs.md}{reproducibility_runs.md}{}.\markdownRendererInterblockSeparator
{}
\markdownRendererSectionEnd \markdownRendererSectionBegin
\markdownRendererHeadingTwo{3. Risultati parabolici}\markdownRendererInterblockSeparator
{}\markdownRendererStrongEmphasis{Bulk e benchmark 1D.} Il caso Okumura, ricondotto a una striscia Q1 uniforme nello spessore con i due estremi dinamici, usa i parametri pubblicati e non sorgenti artificiali. Il confronto tra energia OOFEM e Fig. 4.4 digitalizzata dà RMSE 0.9881, MAE 0.8804 e massimo 1.8323 su 391 campioni nell’intervallo \markdownRendererInlineMath{0.01\le t\le0.4}, con asse energetico circa 0–90. La metrica è stata ricalcolata dai due CSV e coincide col report. È evidenza favorevole per quel problema 1D bulk+capacità/reazione agli estremi, non un confronto con dati numerici tabulati né una verifica del laplaciano tangenziale 2D. La digitalizzazione ha incertezza raster; l’analizzatore energetico usa inoltre Gauss 2×2 sul potenziale quartico Q1, non esatto.\markdownRendererInterblockSeparator
{}\markdownRendererStrongEmphasis{Sistema accoppiato 2D.} La prova parabolica Xiao–Chen usa \markdownRendererInlineMath{\varepsilon=0.1} come ipotesi esplicita e mesh Q1. Il profilo a \markdownRendererInlineMath{T=0.1} è qualitativamente compatibile con la figura, ma il paper non specifica \markdownRendererInlineMath{\varepsilon}, quindi non è una replica quantitativa. A \markdownRendererInlineMath{T=0.1}, gli scarti dalla mesh N=40 sono RMS bulk/traccia .07998/.08963 per N=20 e .06211/.08861 per N=30; la traccia non mostra un miglioramento sostanziale 20→30. In una sequenza a \markdownRendererInlineMath{T=.0625}, dt comune e N=60 come riferimento, il bulk L2 scende .11731→.05461→.02994 per N=20/30/40; la traccia scende .13198→.06707→.03387. Sono scarti interni da riferimento numerico, non errori assoluti.\markdownRendererInterblockSeparator
{}Il raffinamento temporale a N=20 fornisce coppie di differenze RMS mass-weighted bulk 8.6004e-5 e 4.3031e-5 e ordine \markdownRendererInlineMath{p\approx0.999}, coerente con backward Euler. La rianalisi ha rilevato un errore di descrizione nel report: i valori 1.2903e-4 e 4.3031e-5 chiamati “differenze adiacenti in norma di massa” sono in realtà errori rispetto a dt=1e-4; le differenze adiacenti mass-weighted sono quelle riportate sopra. La tendenza d’ordine circa uno resta confermata.\markdownRendererInterblockSeparator
{}\markdownRendererStrongEmphasis{Energia parabolica.} Per \markdownRendererInlineMath{F(u)=\tfrac14(u^2-1)^2}, l’energia potenziale e il bilancio continuo sono\markdownRendererInterblockSeparator
{}\markdownRendererDisplayMath{E_{pot}=\int_\Omega\left(\frac{k}{2}|\nabla u|^2+mF(u)\right)dx+ \int_\Gamma\left(\frac{k}{2}|\nabla_\Gamma\psi|^2+\beta mF(\psi)\right)dS, \qquad \dot E_{pot}=-\int_\Omega c u_t^2dx-\int_\Gamma c\psi_t^2dS\le0.}\markdownRendererInterblockSeparator
{}Nel caso 20×20, \markdownRendererInlineMath{E_0=1.28445597}; a \markdownRendererInlineMath{T=.1}, l’energia è 1.20911974, 1.20901760 e 1.20896702 per dt=.005,.0025,.00125. Il residuo normalizzato cala 4.231e-3→1.201e-3. Nel 16×8, \markdownRendererInlineMath{R(T)} è −2.9468e-4, −1.5248e-4, −7.7664e-5 per dt=.002,.001,.0005. Il trend è coerente col raffinamento, ma la dissipazione è ricostruita con differenze secanti e non è quella esatta integrata nel tempo. Nel grafico secante 2D, il pannello superiore confronta \markdownRendererInlineMath{-\Delta E/dt} con dissipazioni integrate per passo: al primo passo circa 4967.6 contro 0.4959. Quel confronto è dimensionalmente incoerente e non va usato come validazione.\markdownRendererInterblockSeparator
{}
\markdownRendererSectionEnd \markdownRendererSectionBegin
\markdownRendererHeadingTwo{4. Risultati iperbolici}\markdownRendererInterblockSeparator
{}\markdownRendererStrongEmphasis{Bulk.} La prova più forte è l’onda viaggiante Nizovtseva, con parametri mappati \markdownRendererInlineMath{\rho=0.2}, \markdownRendererCodeSpan{wavedamping=5}, k=m=1, \markdownRendererInlineMath{T=.2}, dominio \markdownRendererInlineMath{[-12,12]\times[0,1]}, velocità iniziale analitica. Il massimo mismatch della velocità iniziale è riportato 3.61e-16. A N=80/160/320, gli RMS nodali contro Eq. (29) sono 2.9183e-5, 7.3231e-6, 1.8621e-6; gli ordini spaziali ricalcolati sono 1.9946 e 1.9755. A N=320, le differenze RMS temporali adiacenti sono 2.1093e-6, 5.2419e-7, 1.3885e-7, con ordini 2.0086 e 1.9165 nei due confronti riportati. Sono metriche nodali su nodi x unici, non norme integrate FE. Questo PASS è limitato al bulk, a quel profilo e alla troncatura del dominio.\markdownRendererInterblockSeparator
{}Il modo coseno Newmark contro la soluzione esatta Q1 semidiscreta ha RMS 3.4715e-4, 8.6981e-5, 2.1759e-5 e ordine 1.9968; l’oscillatore uniforme confrontato con DOP853 ha errori 7.3064e-6, 1.8064e-6, 5.0638e-7 e ordine 2.0161 sul primo dimezzamento. Questi isolano temporalità/reazione nel bulk, non la superficie.\markdownRendererInterblockSeparator
{}\markdownRendererStrongEmphasis{Sistema accoppiato.} Il raffinamento interno \markdownRendererInlineMath{dt=.002,.001,.0005} dà ordine 1.9504 contro il run più fine. Un confronto più forte, con ODE semidiscreta assemblata indipendentemente, dà a quei passi RMS bulk 4.6937e-4, 4.6967e-4, 4.6975e-4 e RMS superficiale 2.1058e-3, 2.1070e-3, 2.1073e-3. Gli errori non diminuiscono; gli ordini calcolati sono negativi e vicini a zero. La ripetizione operativa del caso restituisce esattamente il campo OOFEM archiviato e gli stessi errori contro ODE. Questa discrepanza è verificata, ma non individua da sola quale implementazione/reference sia responsabile; non consente di concludere che l’accoppiamento sia validato.\markdownRendererInterblockSeparator
{}
\markdownRendererSectionEnd \markdownRendererSectionBegin
\markdownRendererHeadingTwo{5. Bilanci energetici}\markdownRendererInterblockSeparator
{}Nel ramo iperbolico l’energia totale e il bilancio continuo sono\markdownRendererInterblockSeparator
{}\markdownRendererDisplayMath{E_{tot}=E_{pot}+\int_\Omega\frac{\rho}{2}u_t^2dx, \qquad \dot E_{tot}=-\int_\Omega\rho\zeta u_t^2dx-\int_\Gamma c_\Gamma\psi_t^2dS\le0.}\markdownRendererInterblockSeparator
{}La sola energia potenziale non deve essere usata come test di accettazione perché scambia energia con la cinetica. Nel run singolo 16×8, dt=.0005, la cinetica cresce 0→0.00355156, il potenziale cala di 0.00520321 e il totale scende 1.93160616→1.92995451; il residuo relativo finale è circa 1.80e-7. Nel triplo 16×8 i residui finali sono 3.692e-6, 9.156e-7, 3.473e-7. Questi valori usano la velocità Newmark ricostruita dagli spostamenti; \markdownRendererCodeSpan{Velocity} non è esportata, e l’output testuale arrotondato amplifica l’errore della differenza temporale.\markdownRendererInterblockSeparator
{}Il CSV energetico della prova Nizovtseva è aritmeticamente coerente: E0=3.16242072, ET=3.08549010, residuo finale 3.465e-6. Ma \markdownRendererInlineMath{u_t} è stimata con \markdownRendererCodeSpan{np.gradient} e la velocità iniziale è imposta analiticamente; non è un confronto con lo stato interno Newmark. I CSV iperbolici legacy mostrano discrepanze cinetiche: a dt=.0005 la differenza di potenziale è circa −1.4e-8 mentre quella cinetica è 1.8135e-5; il run finer dt=.00025 ricalcolato dai raw ha R(T)=6.87e-7, contro 3.47e-5 nel CSV legacy. Pertanto i residui/aumenti legacy non sono evidenza affidabile di violazione energetica del solver.\markdownRendererInterblockSeparator
{}La monotonia osservata nei run non dimostra dissipazione discreta generale: backward Euler e Newmark sono valutati a passo finito, con quadratura e derivate temporali specifiche; il codice non dimostra un’identità energetica discreta per il sistema nonlineare accoppiato.\markdownRendererInterblockSeparator
{}
\markdownRendererSectionEnd \markdownRendererSectionBegin
\markdownRendererHeadingTwo{6. Convergenza numerica}\markdownRendererInterblockSeparator
{}\markdownRendererStrongEmphasis{Spaziale:} confronti con soluzione analitica in Nizovtseva e kink sono i più forti; gli ordini sono vicini a due. Nei casi 2D accoppiati parabolici, i valori sono differenze da N=40 o N=60: utili come evidenza di trend, non errori assoluti. La traccia a T=.1 è quasi ferma tra N=20 e N=30. Il raffinamento spaziale accoppiato iperbolico non è completo: manca il CSV e il run/riferimento N=32×16.\markdownRendererInterblockSeparator
{}\markdownRendererStrongEmphasis{Temporale:} Newmark bulk contro eigenmodo semidiscreto/ODE dà ordine circa due; il parabolico 2D dà circa uno. Il test accoppiato interno dà circa 1.95 ma non concorda con la verifica ODE indipendente. Alcuni ordini sono ricavati da sole due differenze: sono indicativi, non dimostrano regime asintotico generale.\markdownRendererInterblockSeparator
{}\markdownRendererStrongEmphasis{Riferimenti fini:} N=40, N=60 e il dt più fine sono riferimenti numerici interni. Errori pari a zero sulla stessa mesh di riferimento sono tautologici. Non vanno chiamati soluzioni esatte né confronti indipendenti.\markdownRendererInterblockSeparator
{}
\markdownRendererSectionEnd \markdownRendererSectionBegin
\markdownRendererHeadingTwo{7. Riproducibilità}\markdownRendererInterblockSeparator
{}Tre casi già documentati sono stati ripetuti in directory temporanee nuove, usando input copiati senza modifiche: parabolico 20×20 (19 s), bulk iperbolico Nizovtseva N=160 (12 s) e accoppiato iperbolico 16×8 (1 s). Tutti hanno terminato con zero errori e warning; i campi finali coincidevano con gli output archiviati. I comandi/log dettagliati sono in \markdownRendererLink{reproducibility\markdownRendererUnderscore{}runs.md}{reproducibility\markdownRendererUnderscore{}runs.md}{reproducibility_runs.md}{}.\markdownRendererInterblockSeparator
{}La ripetibilità è operativa per questi tre casi sul build disponibile, non una riproducibilità bitwise da release: il binario proviene da working tree \markdownRendererCodeSpan{dirty}; parte dei confronti usa \markdownRendererCodeSpan{/tmp}, dipendenze non bloccate e PDF esterni; alcuni runner sovrascrivono output. Il benchmark Lee ha raw ma la metrica L2 non ha analizzatore persistente; Xiao–Chen e Wang/Akrivis mancano dei raw intermedi. La matrice delle metriche e i problemi file-by-file sono in \markdownRendererLink{validation\markdownRendererUnderscore{}metrics\markdownRendererUnderscore{}audit.md}{validation\markdownRendererUnderscore{}metrics\markdownRendererUnderscore{}audit.md}{validation_metrics_audit.md}{} e \markdownRendererLink{reproducibility\markdownRendererUnderscore{}audit.md}{reproducibility\markdownRendererUnderscore{}audit.md}{reproducibility_audit.md}{}.\markdownRendererInterblockSeparator
{}
\markdownRendererSectionEnd \markdownRendererSectionBegin
\markdownRendererHeadingTwo{8. Limiti}\markdownRendererInterblockSeparator
{}\markdownRendererOlBeginTight
\markdownRendererOlItemWithNumber{1}Nessuna prova isolata di inerzia superficiale: l’implementazione corrente non la include; il ramo onda ha inerzia bulk e bordo di primo ordine.\markdownRendererOlItemEnd 
\markdownRendererOlItemWithNumber{2}Il sistema parabolico accoppiato 2D non ha un riferimento esatto o un secondo solver indipendente con dati comparabili; il confronto Xiao–Chen è qualitativo e usa epsilon assunto.\markdownRendererOlItemEnd 
\markdownRendererOlItemWithNumber{3}Il sistema iperbolico accoppiato non supera il cross-check semidiscreto indipendente: errore quasi costante/crescente col raffinamento temporale. La causa resta da isolare.\markdownRendererOlItemEnd 
\markdownRendererOlItemWithNumber{4}Energia cinetica iperbolica è post-processata, non letta direttamente da \markdownRendererCodeSpan{waveVelocity}; parte dei CSV legacy non riproduce i raw.\markdownRendererOlItemEnd 
\markdownRendererOlItemWithNumber{5}Il riferimento Okumura è una curva raster; il benchmark è 1D estruso e il suo post-processing legacy usa quadratura non esatta per il potenziale quartico Q1.\markdownRendererOlItemEnd 
\markdownRendererOlItemWithNumber{6}Test energetici, convergenze spaziali e temporali verificano combinazioni specifiche di mesh, quadratura, dati e schema; non coprono geometrie generiche, coefficienti indipendenti di superficie, condizioni di Dirichlet non omogenee o tutte le opzioni Newmark.\markdownRendererOlItemEnd 
\markdownRendererOlItemWithNumber{7}Riproducibilità da ambiente pulito non è garantita: build sporca, path temporanei, raw mancanti e dipendenze non fissate.\markdownRendererOlItemEnd 
\markdownRendererOlEndTight \markdownRendererInterblockSeparator
{}
\markdownRendererSectionEnd \markdownRendererSectionBegin
\markdownRendererHeadingTwo{9. Conclusioni}\markdownRendererInterblockSeparator
{}\markdownRendererStrongEmphasis{Affermazioni supportate:} il mapping materiale e la forma debole descritti nei report sono coerenti con le equazioni Allen–Cahn adottate; il percorso bulk parabolico ha un confronto bibliografico favorevole in un caso 1D con estremi dinamici; il bulk iperbolico mostra accordo quantitativo con una soluzione analitica non costruita e ordini prossimi a due nei test Newmark disponibili; le metriche principali di questi benchmark sono ricalcolabili dai raw archiviati. Nei casi energetici selezionati, le somme dei contributi recenti sono coerenti e i residui decrescono nel raffinamento iniziale.\markdownRendererInterblockSeparator
{}\markdownRendererStrongEmphasis{Affermazioni eccessive:} “l’intero solver Allen–Cahn è validato”; “la dinamica parabolica di bordo è validata per geometrie 2D generiche”; “il sistema iperbolico accoppiato è validato”; “il ramo Newmark è dissipativo per ogni dt”; “la superficie iperbolica inerziale è stata verificata”; oppure “la monotonia dell’energia dimostra la correttezza”. Queste affermazioni non sono sostenute dalle prove disponibili.\markdownRendererInterblockSeparator
{}\markdownRendererStrongEmphasis{Conclusione scientifica prudente:} l’evidenza è forte ma circoscritta per il bulk iperbolico Nizovtseva e per alcuni aspetti del bulk parabolico. La superficie dinamica parabolica è supportata da benchmark 1D e convergenza interna 2D, ma non ha ancora una validazione quantitativa indipendente completa in 2D. L’accoppiamento iperbolico resta non validato rispetto al riferimento semidiscreto indipendente; la superficie iperbolica di secondo ordine non è un’equazione implementata. Il solver nel suo complesso non può quindi essere dichiarato integralmente validato.
\markdownRendererSectionEnd 
\markdownRendererSectionEnd \markdownRendererDocumentEnd